%% ma: L10-B15-0073
%% lop: 10
%% chuong: 6
%% bai: 15
%% dang: 10.15.1
%% loai: TN4
%% muc: TH
%% dap_an: "A"
%% nguon: "Ảnh giáo viên gửi 25/11/2026: Bài tập trắc nghiệm – Toán 10 Bài 15 Hàm số (3 phần), Phần I Câu 20"
%% kien_thuc: "Lớp 10 Bài 15 (tập xác định của hàm số)"
%% trang_thai: chinh-thuc
%% ghi_chu: "Phân vân TH/VD, chọn TH theo quy tắc mức thấp hơn."
\begin{ex}
Tập xác định của hàm số $y=f(x)=\begin{cases}\sqrt{-3x+8}+x & \text{khi } x<2\\ \sqrt{x+7}+1 & \text{khi } x\ge2\end{cases}$ là
\choice{\True $\mathbb R$}{$\mathbb R\setminus\{2\}$}{$\left(-\infty;\dfrac83\right]$}{$[-7;+\infty)$}
\loigiai{Với $x<2$: $-3x+8>2>0$ nên luôn xác định. Với $x\ge2$: $x+7>0$ nên luôn xác định. Vậy $D=\mathbb R$.}
\end{ex}
